\documentclass[11pt,a4paper]{article}

\def\courseTitle{Industrial Security (Bachelor)}
\def\sectionNumber{08}
\def\sectionTitle{Defense in Depth}
\def\sheetKind{Solution Sheet}

% =====================================================================
% Shared preamble for exercise sheets and solution sheets.
%
% Define BEFORE \input{}-ing this file:
%   \def\courseTitle{...}
%   \def\sectionNumber{...}
%   \def\sectionTitle{...}
%   \def\sheetKind{Exercise Sheet}   or  Solution Sheet
%
% Optional:
%   \def\courseAuthor{...}
%
% The caller must issue \documentclass{...} (typically article) BEFORE
% loading this preamble.
% =====================================================================

\usepackage[utf8]{inputenc}
\usepackage[T1]{fontenc}
\usepackage[english]{babel}
\usepackage[a4paper, margin=2.2cm]{geometry}
\usepackage{graphicx}
\usepackage{listings}
\usepackage{xcolor}
\usepackage{tikz}
\usepackage{booktabs}
\usepackage{array}
\usepackage{enumitem}
\usepackage{titlesec}
\usepackage{fancyhdr}
\usepackage{hyperref}
\usepackage{amsmath}
\usepackage{amssymb}
\usepackage{tcolorbox}
\tcbuselibrary{skins, breakable}
\usepackage{needspace}

% Discourage solo lines split across pages.
\widowpenalty=10000
\clubpenalty=10000
\raggedbottom

% =====================================================================
% Shared TikZ styles for the Industrial Security lecture series.
% Loaded by every slide deck and exercise sheet that uses TikZ.
% =====================================================================

\usetikzlibrary{
  arrows.meta,
  shapes,
  shapes.geometric,
  shapes.symbols,
  shapes.multipart,
  positioning,
  fit,
  calc,
  backgrounds,
  decorations.pathmorphing,
  decorations.pathreplacing,
  decorations.text,
  patterns,
  matrix,
  shadows
}

% --- colour palette --------------------------------------------------
\definecolor{isOT}{HTML}{1F4E79}     % OT (deep blue)
\definecolor{isIT}{HTML}{6F2DA8}     % IT (purple)
\definecolor{isSafety}{HTML}{C00000} % safety red
\definecolor{isNet}{HTML}{2E7D32}    % network green
\definecolor{isAttack}{HTML}{B23A48} % attack red
\definecolor{isAccent}{HTML}{E08E0B} % amber
\definecolor{isMute}{HTML}{6B6B6B}   % grey
\definecolor{isLight}{HTML}{EDF2F8}  % light background
\definecolor{isPLC}{HTML}{2C5282}
\definecolor{isHMI}{HTML}{2F855A}
\definecolor{isSCADA}{HTML}{6B46C1}

% --- generic node styles --------------------------------------------
\tikzset{
  isnode/.style={
    draw, rounded corners=2pt, minimum height=8mm, minimum width=18mm,
    font=\small, align=center, thick
  },
  isbox/.style={isnode, fill=isLight},
  plc/.style={isnode, fill=isPLC!15, draw=isPLC, thick},
  hmi/.style={isnode, fill=isHMI!15, draw=isHMI, thick},
  scada/.style={isnode, fill=isSCADA!15, draw=isSCADA, thick},
  sensor/.style={isnode, fill=isNet!10, draw=isNet, minimum height=6mm, minimum width=14mm, font=\scriptsize},
  actuator/.style={isnode, fill=isAccent!15, draw=isAccent, minimum height=6mm, minimum width=14mm, font=\scriptsize},
  firewall/.style={isnode, fill=isAttack!15, draw=isAttack, thick},
  server/.style={isnode, fill=isMute!15, draw=isMute!80, thick},
  switch/.style={isnode, fill=isLight, draw=black!60, minimum height=6mm, minimum width=14mm, font=\scriptsize},
  workstation/.style={isnode, fill=white, draw=isOT, thick},
  attacker/.style={isnode, fill=isAttack!20, draw=isAttack, thick, font=\small\bfseries},
  inetnode/.style={draw, ellipse, minimum width=24mm, minimum height=10mm, font=\scriptsize, fill=isLight, thick},
  process/.style={isnode, fill=isOT!15, draw=isOT, thick, minimum width=24mm},
  zone/.style={
    draw, rounded corners=4pt, thick, dashed, inner sep=4mm
  },
  conduit/.style={-{Stealth[length=2.5mm]}, thick, isNet!80!black},
  attack/.style={-{Stealth[length=2.5mm]}, thick, isAttack, dashed},
  flow/.style={-{Stealth[length=2.5mm]}, thick, isOT!70!black},
  netline/.style={thick, black!60},
  axisline/.style={-{Stealth[length=2mm]}, thick, black!70},
  tag/.style={font=\scriptsize\itshape, fill=white, inner sep=1pt},
  level/.style={
    draw, fill=isLight, minimum width=0.85\linewidth, minimum height=8mm,
    font=\small, anchor=west
  },
  brace/.style={decorate, decoration={brace, amplitude=4pt, raise=2pt}},
  legendbox/.style={draw, rounded corners=2pt, fill=isLight, inner sep=2mm, font=\scriptsize, align=left}
}

% --- handy macros ----------------------------------------------------
% \plcicon, \hmiicon etc as simple labelled boxes; useful inline.
\newcommand{\plcicon}[1][PLC]{\tikz[baseline=-0.6ex]{\node[plc, minimum width=10mm, minimum height=5mm, font=\scriptsize, inner sep=1pt]{#1};}}
\newcommand{\hmiicon}[1][HMI]{\tikz[baseline=-0.6ex]{\node[hmi, minimum width=10mm, minimum height=5mm, font=\scriptsize, inner sep=1pt]{#1};}}
\newcommand{\fwicon}[1][FW]{\tikz[baseline=-0.6ex]{\node[firewall, minimum width=10mm, minimum height=5mm, font=\scriptsize, inner sep=1pt]{#1};}}


\providecommand{\courseAuthor}{Industrial Security Lecture Series}
\providecommand{\courseInstitute}{Department of Computer Science}
\providecommand{\companyLogo}{logo}

% --- Corporate branding overrides ------------------------------------
% Picks up logo / author / brand colours from the repo-root branding.tex
% when present; falls back to the defaults above otherwise.
\IfFileExists{../../../branding.tex}{% =====================================================================
% Corporate / institutional branding -- single file to customise the
% look of the entire course (slides, notes, exercises, labs).
%
% HOW TO REBRAND IN UNDER A MINUTE:
%   1. Replace `branding/logo.pdf` with your own logo
%      (PDF preferred; PNG and JPG also work -- name it `logo.<ext>`).
%   2. (Optional) change the two colours below to your brand palette.
%   3. (Optional) override the author/institute lines below.
%   4. Re-run `make` at the repo root. That's it.
%
% Everything in this file has sensible defaults; you can delete a line
% to fall back to the built-in defaults, or comment it out.
%
% This file is loaded automatically by the slides, exercise, and lab
% preambles -- you never need to \input it manually.
% =====================================================================

% --- Logo --------------------------------------------------------------
% Path is resolved from each leaf-build directory; the default points at
% the shared `branding/` folder at the repo root. If you put your logo
% somewhere else, set the full or relative path here (without extension):
\renewcommand{\companyLogo}{../../../branding/logo}

% --- Author / institute (appears on the title page footer) ------------
\renewcommand{\courseAuthor}{}
\renewcommand{\courseInstitute}{}

% --- Brand colours ----------------------------------------------------
% slideAccent      = primary brand colour (title bands, accent rule)
% slideSecondary   = secondary brand colour (section badge, highlight)
% Both use HTML hex codes. Keep enough contrast against white text.
\definecolor{slideAccent}{HTML}{00205B}      % deep navy
\definecolor{slideSecondary}{HTML}{00A0DC}   % light blue accent
}{}

% --- listings -------------------------------------------------------
\definecolor{codebg}{rgb}{0.96,0.96,0.96}
\lstset{
  backgroundcolor=\color{codebg},
  basicstyle=\ttfamily\footnotesize,
  breaklines=true,
  frame=single,
  numbers=left,
  numberstyle=\tiny\color{black!50},
  showstringspaces=false,
  tabsize=2
}

% --- header / footer ------------------------------------------------
\pagestyle{fancy}
\fancyhf{}
\renewcommand{\headrulewidth}{0.4pt}
\renewcommand{\footrulewidth}{0pt}
\lhead{\small \courseTitle}
\rhead{\small Section \sectionNumber{} -- \sheetKind}
\cfoot{\thepage}

% --- task / solution boxes -----------------------------------------
% `before={...}` guards every task/solution box with \needspace so the
% box doesn't start with only one or two lines of breathing room at the
% bottom of a page. If less than ~9 lines remain, LaTeX bumps the box
% to the next page; otherwise it starts here and breaks normally if it
% runs long.
% Boxes are `breakable` so very long solutions don't blow the page,
% but `before={\needspace{14\baselineskip}}` pushes them to a fresh
% page whenever fewer than ~14 lines remain. That covers the vast
% majority of single-question splits in practice.
\newtcolorbox{taskbox}[1]{
  enhanced, breakable,
  colback=isLight, colframe=isOT,
  fonttitle=\bfseries\color{white},
  coltitle=white,
  colbacktitle=isOT,
  title={#1},
  arc=2pt, boxrule=0.6pt,
  left=2mm, right=2mm, top=2mm, bottom=2mm,
  before={\par\vspace{2mm}\needspace{14\baselineskip}\par},
  after={\par\vspace{2mm}}
}

\newtcolorbox{solutionbox}{
  enhanced, breakable,
  colback=isNet!5, colframe=isNet,
  fonttitle=\bfseries\color{white},
  coltitle=white,
  colbacktitle=isNet,
  title={Solution},
  arc=2pt, boxrule=0.6pt,
  left=2mm, right=2mm, top=2mm, bottom=2mm,
  before={\par\vspace{2mm}\needspace{14\baselineskip}\par},
  after={\par\vspace{2mm}}
}

\newtcolorbox{hintbox}{
  enhanced, breakable,
  colback=isAccent!10, colframe=isAccent,
  fonttitle=\bfseries\color{white}, coltitle=white, colbacktitle=isAccent,
  title={Hint},
  arc=2pt, boxrule=0.6pt
}

\newtcolorbox{warnbox}{
  enhanced, breakable,
  colback=isAttack!8, colframe=isAttack,
  fonttitle=\bfseries\color{white}, coltitle=white, colbacktitle=isAttack,
  title={Caution},
  arc=2pt, boxrule=0.6pt
}

% --- section formatting --------------------------------------------
\titleformat{\section}{\Large\bfseries\color{isOT}}{\thesection}{0.5em}{}
\titleformat{\subsection}{\large\bfseries\color{isOT!80!black}}{\thesubsection}{0.5em}{}

% --- title block macro ----------------------------------------------
\newcommand{\sheetheader}{%
  \begin{center}
    {\Large\bfseries \courseTitle}\\[0.3em]
    {\large \sheetKind{} -- Section \sectionNumber}\\[0.2em]
    {\large \sectionTitle}\\[0.2em]
    \rule{\linewidth}{0.4pt}
  \end{center}
  \vspace{0.5em}
}


\begin{document}
\sheetheader

\noindent\textbf{Rubric note.} These are reference answers, not the only acceptable ones. Award credit for any answer that (a) names a defensible architecture, (b) cites a real product or standard, and (c) acknowledges the operational trade-off. Penalise flat-VPN designs, ``patch everything'' answers, and any solution that hand-waves the audit trail.

\vspace{0.6em}

\begin{solutionbox}
\textbf{Exercise 8.1 -- Vendor PLC Download Workflow.}
Path: Siemens engineer laptop $\to$ public Internet $\to$ brokered portal in the iDMZ (Cyolo, Claroty SRA, or BeyondTrust Privileged Remote Access) with TOTP or FIDO2 MFA $\to$ Windows jump host in iDMZ running TIA Portal, with full keystroke and screen recording (BeyondTrust session manager or Cyolo recording) $\to$ industrial firewall (Tofino, Cisco ISA-3000, Fortinet Rugged) enforcing an S7Comm-only ACL with time-of-day restriction $\to$ S7-1500. Operator approval is a click in the broker by the on-shift control-room operator; the request expires after a 30--60 minute window. Audit trail: broker auth log, MFA event, operator approval click, full session video, jump-host command history, firewall flow log, PLC project checksum before/after. Failure mode -- broker outage: design ships with a documented break-glass procedure (physical key to the cabinet, two-person rule, post-hoc audit within 24 h); never a permanent VPN as fallback.
\textit{Rubric: must name broker + jump host + MFA + operator approval + recording, and one named product. Half credit for omitting the audit-trail enumeration.}
\end{solutionbox}

\begin{solutionbox}
\textbf{Exercise 8.2 -- Counter-proposal to Flat VPN.}
Acknowledge the legitimate need: quarterly firmware drops, vibration tuning, warranty diagnostics. Propose a brokered remote-access service (Claroty SRA, Cyolo, BeyondTrust PRA) terminating in the iDMZ, with operator-approved sessions, full recording, and per-action MFA. Map to IEC 62443-3-3 SR\,1.13 (access via untrusted networks) and SR\,2.4 (mobile code). The three unacceptable risks of flat VPN: (i) lateral movement -- once the tunnel is up the vendor laptop is on the plant LAN with no per-host control; (ii) persistence -- a compromised vendor laptop keeps the tunnel alive 24/7; (iii) audit gap -- you cannot tell which engineer did what, only that ``a Siemens IP wrote to the PLC''. Pilot: 30 days, one turbine, vendor keeps their existing change-window cadence, broker covers their TIA Portal and OPC traffic only. Offer to fund the vendor account in the broker so cost is not their objection.
\textit{Rubric: must offer an alternative, not just refuse. Credit for citing 62443-3-3 or NIST SP 800-46.}
\end{solutionbox}

\begin{solutionbox}
\textbf{Exercise 8.3 -- Passive OT IDS Placement.}
One sensor per substation (3) at the managed switch SPAN port, mirroring the process VLAN (IEC 61850 GOOSE/MMS, DNP3, IEC 60870-5-104) but \emph{not} the management VLAN. One sensor at the central site on a hardware TAP between the SCADA HMI and the historian -- TAP rather than SPAN because the SCADA polling load risks SPAN drops on a busy switch. Total: 4 sensors, central management plane (Malcolm/Nozomi CMC/Claroty EMC) on the iDMZ. What you miss: encrypted vendor management tunnels into the IEDs, anything that traverses the substation cellular backhaul without crossing the SPAN port, and any traffic east-west between two IEDs on the same VLAN if your switch SPAN only mirrors uplinks. Compensating controls: enforce vendor tunnel termination through the broker; deploy host-based logging on the IEDs that support it; configure SPAN to include the inter-IED traffic explicitly.
\textit{Rubric: must distinguish SPAN vs TAP and name one blind spot. Credit for explicitly stating which VLAN is mirrored.}
\end{solutionbox}

\begin{solutionbox}
\textbf{Exercise 8.4 -- Three Layers, One Logo.}
Example: perimeter firewall, internal firewall, and host firewall all from the same vendor, all managed by the same management console with the same admin credentials. A single credential theft or a single zero-day in the management plane defeats all three layers simultaneously. Equivalent failure modes: three layers of detection all dependent on the same Active Directory; three AV products all using the same signature feed; three backups all written from the same backup agent with the same service account. Restore independence by mixing vendors at one layer (e.g.\ a Fortinet plant firewall behind a Cisco perimeter), using an out-of-band identity for the most critical management plane (a separate Tier-0 AD or local accounts), or running the management plane on a different OS family (Linux jump host fronting a Windows admin tier).
\textit{Rubric: must name the shared dependency, not just ``same vendor''. Credit for an out-of-band identity proposal.}
\end{solutionbox}

\begin{solutionbox}
\textbf{Exercise 8.5 -- ``Just 5 Minutes With My Own Laptop''.}
Default answer: no. Backing: IEC 62443-3-3 SR\,1.1 (identification and authentication of human users) and SR\,1.2 (software process and device identification), reinforced by the plant's BYOD-prohibition policy and INL DiD 2016 \S 5 (physical and network access control). Workaround when the read cannot wait: contractor copies the diagnostic file to a USB stick, kiosk (OPSWAT MetaDefender) scans and re-emits onto a sanitised stick, an authorised engineer mounts the sanitised stick on a hardened plant workstation. If the file must come from the laptop directly, terminate the laptop in a vendor DMZ VLAN (NAC quarantine class) and broker a one-way file transfer via a jump host. Pre-conditions for any ``yes'': NAC profile assigns vendor-VLAN by certificate or one-time RADIUS PSK, EDR scan acceptable, time-boxed (30 min) authorisation ticket open, on-shift operator approves, session recorded. Audit-trail entries: ticket number, contractor identity, MAC observed by NAC, VLAN assigned, files transferred (with hashes), supervisor approval, ticket closure.
\textit{Rubric: must default to ``no'' before offering an alternative; must place the laptop somewhere other than the OT VLAN.}
\end{solutionbox}

\begin{solutionbox}
\textbf{Exercise 8.6 -- Backup Plan for a Mixed Plant.}
\begin{itemize}\setlength{\itemsep}{1pt}\small
  \item \textbf{M580 PLC} -- EcoStruxure Control Expert project + checksum, on change, 5 versions then quarterly snapshot, off-site immutable, annual restore-test on FAT bench.
  \item \textbf{Wonderware HMI} -- AVEVA System Platform Galaxy backup + Veeam VM image, daily, 30\,d incremental + 12\,mo full, off-site immutable, semi-annual restore.
  \item \textbf{SQL historian} -- native SQL full + diff + log plus Veeam VM, full weekly / diff daily / log 15\,min, 90\,d online + 12\,mo archive, immutable bucket (S3 Object Lock), quarterly point-in-time restore.
  \item \textbf{12 managed switches} -- RANCID or Oxidized running-config pull, on change + daily, 1\,year retention, off-site, annual rebuild-of-spare test.
\end{itemize}
\noindent 3-2-1 holds: two media (disk + immutable object store), at least one off-site, the immutable copy is the ransomware insurance.
\textit{Rubric: must give a frequency, retention, off-site, and restore-test for each asset. Credit for immutable storage; deduct for ``daily full of 300 GB historian''.}
\end{solutionbox}

\begin{solutionbox}
\textbf{Exercise 8.7 -- Compensating Controls for Windows 7 HMI.}
Priority order: (1) micro-segmentation -- put the HMI behind a Tofino or Cisco ISA-3000 with an allow-list of only the PLC IPs and protocols it actually needs; protects against lateral movement from compromised IT. (2) Application allow-listing -- McAfee Solidcore or AppLocker in enforce mode pinned to the frozen HMI binary set; protects against arbitrary code execution. (3) Virtual patching -- Trend Micro TippingPoint or Snort/Suricata signatures in front of the HMI covering the unpatched CVEs; protects against known network-borne exploits. (4) Disable USB at OS layer via Group Policy plus McAfee Device Control; protects against the malicious-USB vector. (5) Enhanced monitoring -- forward all Windows event logs and any HMI alarm exports to the OT-aware SIEM with tight alerts on logon, service creation, and registry persistence; protects nothing but cuts dwell time. Residual risk: a determined attacker with valid operator credentials and physical access to the console can still abuse the application allow-listed HMI itself -- no compensating control replaces a working patch for the underlying OS.
\textit{Rubric: must rank in priority order with one named product per control. Credit for an honest residual-risk statement.}
\end{solutionbox}

\begin{solutionbox}
\textbf{Exercise 8.8 -- OT Log Sources Splunk Misses.}
(1) Cisco Cyber Vision sensor alerts (S7Comm/Modbus session anomalies) -- onboard via the Cyber Vision Splunk add-on, CIM map to \texttt{Network\_Traffic} and \texttt{Alerts}. (2) IEC 60870-5-104 ASDU events from a Zeek dissector -- ship Zeek logs through Cribl or the Splunk add-on for Zeek, custom CIM mapping for control-direction commands. (3) PLC RUN/STOP transitions from a Claroty CTD or Dragos Platform feed -- syslog forwarder into Splunk with a Dragos technology add-on. (4) Wonderware/AVEVA Galaxy alarm-history rows -- DB Connect or ODBC scheduled query, normalise to \texttt{Change} data model. (5) BeyondTrust PRA / Cyolo session metadata (session-start, target host, recording URL) -- vendor Splunk app or CEF over syslog, normalise to \texttt{Authentication} and \texttt{Change}. Cross-source correlation rule: a successful jump-host session by vendor X (source 5), followed within 10 minutes by a PLC RUN$\to$STOP transition (source 3) on a PLC that source 1 attributes to that jump host. Generic Splunk ES will not write that rule for you -- the OT add-on is what makes it possible.
\textit{Rubric: must name five real OT sources and a real ingestion path. Credit for a correlation rule that crosses at least two of them.}
\end{solutionbox}

\begin{solutionbox}
\textbf{Exercise 8.9 -- Honey-PLC Deployment Plan.}
Product: Thinkst Canary for a turnkey, supported deployment, with conpot as an open-source secondary on a Raspberry Pi for breadth. Exposes Modbus/TCP 502, S7Comm 102, EtherNet/IP 44818, and a fake Siemens Webserver on 80/443. Naming convention matches the real fleet: \texttt{LINE3-PLC-07}, \texttt{LINE4-PLC-11}, etc., so an attacker enumerating asset names cannot tell the decoy from the real PLC. Placement: one per process VLAN at Purdue Level 2 (3--5 total), one in the iDMZ at Level 3.5, one on the engineering workstation VLAN. Alerts: any TCP SYN to any honey-PLC port writes a high-severity event to the SIEM, posts to the SOC chat channel, and pages the on-call. Triage: because legitimate traffic to a honey-PLC is zero by design, a single hit is treated as a confirmed incident -- isolate the source, capture pcap, kick off the IR playbook.
\textit{Rubric: must name a product, expose realistic OT protocols, and define how an alert reaches a human. Credit for the ``zero false positive'' framing.}
\end{solutionbox}

\begin{solutionbox}
\textbf{Exercise 8.10 -- Spot the Bug: NAC Config.}
Three weaknesses. (i) MAC OUI is trivially spoofable -- any attacker can set their laptop NIC to a Dell OUI in under a minute. (ii) ``Dell Inc.'' is one of the largest OUI blocks in existence; the rule effectively admits half the laptops on the planet. (iii) Level 3 supervisory is the wrong VLAN for a transient IT-vendor device -- it places the laptop adjacent to engineering workstations and historians. Three fixes, increasing cost. (a) Replace MAB with EAP-TLS using a per-device certificate issued by the plant CA -- supported by IEEE 802.1X-2020. (b) Add a posture check (Forescout SilentDefense or Cisco ISE) that verifies EDR present, OS patch level, and disk encryption before granting any VLAN -- aligns with NIST SP 800-82 R3 \S 6 access-control guidance. (c) Terminate the vendor laptop in a vendor DMZ VLAN with no direct Level 3 access; required action goes through a brokered jump host -- aligns with IEC 62443-3-3 SR\,1.1 and the zone/conduit model.
\textit{Rubric: must catch the MAC-spoofing weakness explicitly and propose at least one certificate-based or brokered alternative.}
\end{solutionbox}

\end{document}
